\documentclass[10pt,a4paper]{amsart}
\usepackage{amsmath,amssymb,amsthm,graphicx}
\usepackage{verbatim}
\newenvironment{reference}{\comment}{\endcomment}
%\newenvironment{reference}{}{}
\newenvironment{slogan}{\comment}{\endcomment}
\newenvironment{history}{\comment}{\endcomment}
% For commutative diagrams we use Xy-pic
\usepackage[all]{xy}

% We use 2cell for 2-commutative diagrams.
\xyoption{2cell}
% Theorem environments.
%
\theoremstyle{plain}
\newtheorem{theorem}[subsection]{Theorem}
\newtheorem{proposition}[subsection]{Proposition}
\newtheorem{lemma}[subsection]{Lemma}

\theoremstyle{definition}
\newtheorem{definition}[subsection]{Definition}
\newtheorem{example}[subsection]{Example}
\newtheorem{exercise}[subsection]{Exercise}
\newtheorem{situation}[subsection]{Situation}

\theoremstyle{remark}
\newtheorem{remark}[subsection]{Remark}
\newtheorem{remarks}[subsection]{Remarks}

\numberwithin{equation}{subsection}

% Macros
%
\def\lim{\mathop{\mathrm{lim}}\nolimits}
\def\colim{\mathop{\mathrm{colim}}\nolimits}
\def\Spec{\mathop{\mathrm{Spec}}}
\def\Hom{\mathop{\mathrm{Hom}}\nolimits}
\def\Ext{\mathop{\mathrm{Ext}}\nolimits}
\def\SheafHom{\mathop{\mathcal{H}\!\mathit{om}}\nolimits}
\def\SheafExt{\mathop{\mathcal{E}\!\mathit{xt}}\nolimits}
\def\Sch{\mathit{Sch}}
\def\Mor{\mathop{\mathrm{Mor}}\nolimits}
\def\Ob{\mathop{\mathrm{Ob}}\nolimits}
\def\Sh{\mathop{\mathit{Sh}}\nolimits}
\def\NL{\mathop{N\!L}\nolimits}
\def\CH{\mathop{\mathrm{CH}}\nolimits}
\def\proetale{{pro\text{-}\acute{e}tale}}
\def\etale{{\acute{e}tale}}
\def\QCoh{\mathit{QCoh}}
\def\Ker{\mathop{\mathrm{Ker}}}
\def\Im{\mathop{\mathrm{Im}}}
\def\Coker{\mathop{\mathrm{Coker}}}
\def\Coim{\mathop{\mathrm{Coim}}}

% Boxtimes
%
\DeclareMathSymbol{\boxtimes}{\mathbin}{AMSa}{"02}

%
% Macros for moduli stacks/spaces
%
\def\QCohstack{\mathcal{QC}\!\mathit{oh}}
\def\Cohstack{\mathcal{C}\!\mathit{oh}}
\def\Spacesstack{\mathcal{S}\!\mathit{paces}}
\def\Quotfunctor{\mathrm{Quot}}
\def\Hilbfunctor{\mathrm{Hilb}}
\def\Curvesstack{\mathcal{C}\!\mathit{urves}}
\def\Polarizedstack{\mathcal{P}\!\mathit{olarized}}
\def\Complexesstack{\mathcal{C}\!\mathit{omplexes}}
% \Pic is the operator that assigns to X its picard group, usage \Pic(X)
% \Picardstack_{X/B} denotes the Picard stack of X over B
% \Picardfunctor_{X/B} denotes the Picard functor of X over B
\def\Pic{\mathop{\mathrm{Pic}}\nolimits}
\def\Picardstack{\mathcal{P}\!\mathit{ic}}
\def\Picardfunctor{\mathrm{Pic}}
\def\Deformationcategory{\mathcal{D}\!\mathit{ef}}
\usepackage{cleveref}
\crefname{originallemma}{original lemma}{original lemmas}
\Crefname{originallemma}{Original lemma}{Original lemmas}
\makeatletter
\newlabel{lemma-intersection-closed-in-quasi-compact}{{lemma-intersection-closed-in-quasi-compact}{source}{}{sourcepointer-1}{}}
\newlabel{lemma-intersection-closed-in-quasi-compact@cref}{{[originallemma][0][]lemma-intersection-closed-in-quasi-compact}{[1][0][]source}{}{}{}}
\makeatother
\crefname{originallemma}{original lemma}{original lemmas}
\Crefname{originallemma}{Original lemma}{Original lemmas}
\makeatletter
\newlabel{lemma-inverse-limit-quasi-compact}{{lemma-inverse-limit-quasi-compact}{source}{}{sourcepointer-2}{}}
\newlabel{lemma-inverse-limit-quasi-compact@cref}{{[originallemma][0][]lemma-inverse-limit-quasi-compact}{[1][0][]source}{}{}{}}
\makeatother
\title{An inverse limit of nonempty quasi-compact Hausdorff spaces is nonempty}
\author{The Stacks Project Authors}
\date{stacks-project:topology.tex:local-snapshot-20261012; GNU FDL 1.2 or later}
\begin{document}
\maketitle
\noindent\textit{Editorial note.} All delivered mathematical text is the author's own, copied from the pinned Stacks Project chapter source sources/stacks-topology/pinned-original-source.tex; no statement or proof step was rewritten, reordered or removed.\par\medskip
\noindent\textit{Editorial note.} Preamble provenance: the Stacks Project keeps one single shared preamble for every chapter; this card takes it directly from the staged sources/stacks-topology/preamble.tex as three verbatim spans (lines 8-12, 14-18 and 154-217), so no derived file was written and nothing was added or redefined. The wrapper class and the base packages are supplied by the assembler.\par\medskip
\noindent\textit{Editorial note.} Bibliography provenance: neither the counted statement nor its proof contains a citation (this lane rescans the whole span before writing the card), so the card carries no bibliography and no entry was reconstructed.\par\medskip
\noindent\textit{Editorial note.} Selection provenance: this card was opened under the file-based control contract (CONTROL.md), with the mechanical pre-card gate run on its statement and proof interval before the card was written.\par\medskip
\noindent\textit{Allowed standard prerequisites (editorial).}\begin{itemize}
\item The description of the limit as the intersection of the closed subsets cut out by the graphs of the transition maps inside the product. Invoked at the start of the counted proof (source lines 2732-2740). Presents the limit as an intersection of closed subsets of a quasi-compact space.
\item The statement that an intersection of closed subsets of a quasi-compact space is nonempty when every finite subfamily has nonempty intersection. Invoked inside the counted proof (source lines 2737-2739). Reduces nonemptiness of the limit to the finite intersection property.
\end{itemize}\medskip
\ifdefined\hypertarget\hypertarget{sourcepointer-1}{}\fi
\noindent\textit{Source reference.} Original lemma lemma-intersection-closed-in-quasi-compact: sources/stacks-topology/pinned-original-source.tex, source lines 1941--1941. Navigation pointer only: the counted text refers to the Stacks item lemma-intersection-closed-in-quasi-compact, which is an existing labeled result of this same chapter and is not part of this card. The pointer carries the target's own label as its number and its definition line as its span; its content is not reproduced here.\par\medskip
\ifdefined\hypertarget\hypertarget{sourcepointer-2}{}\fi
\noindent\textit{Source reference.} Original lemma lemma-inverse-limit-quasi-compact: sources/stacks-topology/pinned-original-source.tex, source lines 2692--2692. Navigation pointer only: the counted text refers to the Stacks item lemma-inverse-limit-quasi-compact, which is an existing labeled result of this same chapter and is not part of this card. The pointer carries the target's own label as its number and its definition line as its span; its content is not reproduced here.\par\medskip
\begin{lemma}
\label{lemma-nonempty-limit}
Let $\mathcal{I}$ be a cofiltered category and let $i \mapsto X_i$
be a diagram over $\mathcal{I}$ in the category of topological
spaces. If each $X_i$ is quasi-compact, Hausdorff, and nonempty, then
$\lim X_i$ is nonempty.
\end{lemma}

\begin{proof}
In the proof of Lemma \ref{lemma-inverse-limit-quasi-compact}
we have seen that $X = \lim X_i$ is the intersection of the
closed subsets
$$
Z_\varphi = \Gamma_\varphi \times \prod\nolimits_{l \not = j, k} X_l
$$
inside the quasi-compact space $\prod X_i$ where $\varphi : j \to k$
is a morphism of $\mathcal{I}$ and $\Gamma_\varphi \subset X_j \times X_k$
is the graph of the corresponding morphism $X_j \to X_k$. Hence by
Lemma \ref{lemma-intersection-closed-in-quasi-compact}
it suffices to show any finite intersection of these subsets is nonempty.
Assume $\varphi_t : j_t \to k_t$, $t = 1, \ldots, n$ is a finite collection
of morphisms of $\mathcal{I}$. Since $\mathcal{I}$ is cofiltered, we can
pick an object $j$ and a morphism $\psi_t : j \to j_t$ for each $t$.
For each pair $t, t'$ such that either (a) $j_t = j_{t'}$, or
(b) $j_t = k_{t'}$, or (c) $k_t = k_{t'}$ we obtain two morphisms
$j \to l$ with $l = j_t$ in case (a), (b) or $l = k_t$ in case (c).
Because $\mathcal{I}$ is cofiltered and since there are finitely
many pairs $(t, t')$ we may choose a map $j' \to j$ which equalizes
these two morphisms for all such pairs $(t, t')$. Pick an element
$x \in X_{j'}$ and for each $t$ let $x_{j_t}$, resp.\ $x_{k_t}$
be the image of $x$ under the morphism $X_{j'} \to X_j \to X_{j_t}$,
resp.\ $X_{j'} \to X_j \to X_{j_t} \to X_{k_t}$.
For any index $l \in \Ob(\mathcal{I})$ which is not equal to
$j_t$ or $k_t$ for some $t$ we pick an arbitrary element $x_l \in X_l$
(using the axiom of choice). Then $(x_i)_{i \in \Ob(\mathcal{I})}$
is in the intersection 
$$
Z_{\varphi_1} \cap \ldots \cap Z_{\varphi_n}
$$
by construction and the proof is complete.
\end{proof}

\end{document}
